% Pending research note. Not an included result and not a reduction of A5's gate.
\subsection*{Two separate failures in a Thom-recognition approach}

The first statement concerns existence of any finite Thom stratification.
The second concerns using a fixed smooth fibre as a source stratum.
They must not be conflated.

\medskip\noindent\textbf{1. A proper real algebraic map with no finite
Thom \(a_f\) stratification.}
Let
\[
 X=S^2=\{(x,t,z)\in\mathbb R^3:x^2+t^2+z^2=1\},\qquad
 b:X\longrightarrow\mathbb R^2,\qquad b(x,t,z)=(x,xt).
\]
This is a polynomial map on a compact smooth real algebraic manifold;
in particular it is proper. Its fibre at the origin is the circle
\(E=\{x=0\}\subset S^2\).

For clarity, a finite \(C^1\) semialgebraic Thom stratification here means
finite source and target stratifications, with the frontier condition,
such that every source stratum maps submersively to a target stratum
and every incident source pair \(S,T\) satisfies \(a_b\):
if \(q_i\in S\) tends to \(q\in T\) and the tangent spaces to the
fibres of \(b|_S\) tend to a plane \(L\), then
\[
                    \ker d(b|_T)_q\subset L.
\]
The analogous definable-curve formulation works over a real closed field.

No such finite stratification exists. On the open subset \(xz\ne0\),
the map has rank two: in the chart \(z\ne0\), its Jacobian determinant
in coordinates \((x,t)\) is \(x\). Equivalently the determinant of
the two coordinate gradients of \(b\) together with the sphere
gradient is \(2xz\). At every point of \(E\) its rank is one.
Thus a two-dimensional source stratum cannot contain a point of \(E\):
it would be open in the sphere near that point, would contain nearby
rank-two points, and could not map submersively to one fixed target
stratum of a single dimension.

A finite stratification of the circle by source strata of dimension
at most one has a one-dimensional stratum \(T\) containing an open
arc of \(E\); choose an interior point \(q\) of such an arc with
\(z(q)\ne0\). The intersection with \(E\) is open in \(T\) near \(q\),
so \(d(b|_T)_q=0\) and its kernel is the nonzero tangent line \(T_qE\).
The union of all source strata of dimension less than two has empty
interior in the sphere. Hence points of two-dimensional source strata
with \(xz\ne0\) approach \(q\). Finiteness gives one such stratum \(S\)
with \(q\in\overline S\); the frontier condition gives
\(T\subset\overline S\). Definable curve selection gives an approach
inside \(S\cap\{xz\ne0\}\). Every fibre tangent along that approach is
zero, since \(b|_S\) has rank two. Thom's condition would therefore
require \(T_qE\subset0\), a contradiction.

The local polynomial \((x,t)\mapsto(x,xt)\) is the blowdown obstruction
appearing on printed p.33 of Shiota's original
\emph{O-minimal Hauptvermutung for polyhedra}, arXiv:1002.1508.
The compact-sphere version and the tangent-space proof above are
written out here. This obstruction already has two-dimensional target;
it is not a defect caused only by integrating a nondefinable vector field.
It refutes the proposed universal premise that proper real definable
maps admit finite Thom \(a_f\) stratifications.

\medskip\noindent\textbf{2. A proper trivial bundle can fail \(a_f\)
along its entire smooth special fibre.}
Set
\[
 g(x,t)=(t^2-x)^2+\frac{x^2}{4},\qquad f(x,t)=xg(x,t),
 \qquad
 H(x,t)=\bigl(f(x,t),\,t\sqrt{g(x,t)}\bigr).
\]
The only zero of \(g\) is the origin. With the weighted dilation
\(\delta_s(x,t)=(s^2x,st)\), one has
\(g(\delta_s(x,t))=s^4g(x,t)\). Direct substitution gives
\[
                    g(H(x,t))=g(x,t)^3.
\]
Consequently \(H:\mathbb R^2\to\mathbb R^2\) has the explicit inverse
\[
 H^{-1}(y,v)=
 \begin{cases}
 \bigl(y/g(y,v)^{1/3},\,v/g(y,v)^{1/6}\bigr),&(y,v)\ne(0,0),\\
 (0,0),&(y,v)=(0,0).
 \end{cases}
\]
All roots are the positive real roots. Both maps are semialgebraic.
Continuity of the inverse at the origin follows from
\[
 g(y,v)\ge\frac{y^2}{4},\qquad
 g(y,v)\ge\frac{v^4}{5};
\]
the latter is the identity
\[
 g(y,v)-\frac{v^4}{5}
             =\frac54\left(y-\frac45v^2\right)^2.
\]
Writing \(a=g(y,v)\), the two inverse coordinates are bounded by
\(2a^{1/6}\) and \(5^{1/4}a^{1/12}\), respectively.
Thus \(H\) is a semialgebraic homeomorphism.

Let
\[
 X=H^{-1}([-1,1]^2),\qquad f_X:X\longrightarrow[-1,1].
\]
Equivalently \(X\) is defined by \(f^2\le1\) and \(t^2g\le1\).
It is closed and bounded: the inequalities above give
\(x^6\le16\) and \(t^6\le5\). Hence \(f_X\) is proper.
The map \(H|_X\) identifies it with the product projection
\([-1,1]\times[-1,1]\to[-1,1]\). Its maximal open manifold
triviality locus is all of \((-1,1)\).

Nevertheless, along the entire smooth fibre
\(F=f_X^{-1}(0)\), which is an interval on the vertical axis, \(a_f\)
fails at the origin. Indeed,
\[
 f_x=\frac{15}{4}\left(x-\frac25t^2\right)
                       \left(x-\frac23t^2\right),
 \qquad
 f_t=4xt(t^2-x).
\]
At \(x=\frac25t^2\), \(t\ne0\), the first derivative is zero and
the second is \(\frac{24}{25}t^5\ne0\).
The tangent to the nearby fibre is therefore horizontal,
whereas \(T_0F\) is vertical.

Even the following weaker retraction condition fails for every
\(C^1\) local retraction \(\pi\) onto \(F\):
\[
       d\pi_q:\ker df_q\longrightarrow T_{\pi(q)}F
                    \quad\hbox{is surjective for all nearby }q\notin F.
\]
Write \(\pi(x,t)=(0,\phi(x,t))\), so
\(\phi(0,t)=t\), \(\phi_t(0,0)=1\), and \(\phi_x\) is locally bounded.
On the curve \(x=kt^2\), divide the derivative of \(\phi\) on the
fibre-tangent vector \((f_t,-f_x)\) by \(t^4\). As \(t\to0\) this tends to
\[
             -P(k),\qquad P(k)=1-4k+\frac{15}{4}k^2.
\]
The values of \(P\) at \(1/4,1/2,1\) are respectively
\(15/64,-1/16,3/4\).
For each sufficiently small positive \(t\), continuity in \(k\)
therefore gives a zero between \(1/4\) and \(1/2\), and another
between \(1/2\) and \(1\).
The gradient of \(f\) vanishes only at the origin: the alternatives
from \(f_t=0\) are \(x=0\), \(t=0\), or \(x=t^2\), and direct
substitution in \(f_x\) proves the claim.
At each of the obtained points the nonzero vector
\((f_t,-f_x)\) spans \(\ker df\), and \(d\pi\) annihilates it.
This proves failure of the weaker condition arbitrarily near the origin.

This odd-degree modification is motivated by Koike--Shiota,
\emph{Non-smooth points set of fibres of a semialgebraic mapping},
J. Math. Soc. Japan 59 (2007), Example 4.7 and Remark 4.8,
printed pp.966--968. Their original uses
\(x^2((t^2-x)^2+x^2/20)\) and studies triviality along the fibre;
the displayed \(f=xg\) here is a full product over the target.
The original article was opened in full at
\url{https://www.academia.edu/21193908/Non_smooth_points_set_of_fibres_of_a_semialgebraic_mapping}.
No assertion that this second example rules out every finite
Thom stratification is made: splitting \(F\) at the origin
removes that specific obstruction.

\medskip\noindent\textbf{Exact algebra certificate.}
The program \path{code/check_thom_obstructions.py} verifies the two
Jacobian determinants, the collapse equation, the weighted-homeomorphism
identity, the positivity identities, the derivative factors, and
the three sign values over \(\mathbb Q\).
Its actual output is \path{logs/check_thom_obstructions.txt};
the recorded runtime was \(0.365916400\) seconds.
These finite algebra checks complement the full topological
arguments above; they do not purport to enumerate stratifications.

\medskip\noindent\textbf{Consequence for the remaining A.5 gate.}
The universal finite-Thom-stratification premise is false,
while first-isotopy-lemma loci are only sufficient loci.
Neither example refutes definability of the maximal product locus,
and neither reduces the currently remaining case of fibre dimension
at least two and target dimension at least three.
